\subsection{Realizing both reselected sets on the prescribed windows}
\label{scope:reselected-realization}

A family already chosen has very little freedom at a fixed scale. For any
two of its indices \(i,i'\) at \(x\), the integer encodings of their top
sets both equal \(\mathtt{Fm.TTn}(x)\), and their integer thick sets agree
exactly. These are the fields \code{TT\_eq} and
\code{realisedChain\_torus\_eq} of \code{FamCore}. The map assigning an
index to each translate therefore cannot implement a strict thinning in
that family. The construction below chooses a new family.

\begin{proposition}[Common weights and points on a prescribed host]
\label{scope:reselected-hosts}
Fix \(0<\varepsilon<1/52\) and standard scale constants \(C\).
For \(c_D\ge0\), \(c_h\ge402\) and any fixed host bound \(C_8\), there is
a threshold, uniform in the host, above which every \code{RankVal},
\code{HostVal} host has one pair of integer centre vectors and genuine
Harish--Chandra parameters for every code word at every translate.

For the original geometric point construction take the staged large
parameters of \code{ObjCP}. This gives a positive upper bound on
\(\varepsilon\). After fixing a positive \(\varepsilon\) below that bound,
\(C_8\ge 1\), and \(C_5,c_6,C_b\),
the size requirement on \(C\) is chosen before \(C\). Take standard
\(C\) meeting this requirement with \(C.\mathtt{P.Cpde}=C_8\). Then \(c_D\) is
chosen sufficiently large, \(c_h\) sufficiently large in terms of
\(c_D\), and finally a scale threshold is fixed. At every later scale
with exponent \(\varepsilon\), prescribe a \code{RankVal} host \(H\) satisfying
\(H.C_{T_2}=C.C_{T_2}\) and \code{HostGood} with the canonical auxiliary
level prime \(q_0(H)\). There is one kit of actual points from \(H\) covering
all \(0\le u<n\). Its centres and injective word code are common to all
translates. If such a host is prescribed at every sufficiently late scale
in a window set \(S\), one can choose a family with
\(\mathtt{Fm.Win}=S.\mathtt{Win}\). Any property of the scale, ranks,
raw host with its field axioms, top set and integer thick set that holds
on the prescribed hosts also holds at their actual points. This includes the cohomology bound
\(\mathrm{VVle}(E,R.t,\sqrt b)\) and an overlap estimate on these same
two sets.
\end{proposition}

\begin{proof}
The elementary centre construction \code{exists\_weights\_in\_windows}
takes \(a_t=\lceil a\rceil\), \(b_t=\lceil b'\rceil\) on the top set.
Both windows have width at least one beyond a uniform threshold.
Having fixed these vectors, apply \code{exists\_hc\_family} to each
translate of the given host. For \(0\le j<m\), put
\(h=\lceil x^{c_h}\rceil\). At a given embedding \(\tau\), its two arrays
are \(b_j=B_\tau+2jh\) and \(e_j=B_\tau+(2j+1)h\).
At a top embedding \(B_\tau=b_t+\lfloor b'/x\rfloor\), and at a thin
embedding \(B_\tau=\lceil w''(n)\rceil\). Start every thick embedding at
\(\lceil b\rceil\). At one thick embedding the construction adds an
integer between zero and \(\lvert\mu(E)\rvert-1\) to satisfy the product
condition on the roots of unity of \(E\). The arrays \(b,e,v\) and the
two end entries of \(u\) are shared among the code words at a fixed
translate. Each interior entry of \(u\) records its own code bit. This
proves the weight assertion directly on the selected host.

The original automorphic existence theorem applies to every admissible
host, to every pair of centres in \code{WinOK}, to every injective code of
the required size and to every \(u<n\), in that order. The objects in this
assertion come from the two automorphic constructions in
\code{MemberE}. Its four fields give the base condition, first construction,
export and second construction used by
\code{memberTop\_final\_contextParameters}. More precisely, for each
admissible pre-chain and its centres and code, \code{RecExP} first
chooses the positive word-indexed integers \(d_u\) and a positive real
\(D_*\). For every level choice with these values and the specified
\(\Gamma\), it supplies a single \(D\) and a layer \(Y\) over that level.
The reconstruction, ledger and local quadratic assertions all hold on
this same pair. The final object property applies to every compatible
part of \(Y\). Taking the final object
predicate to be true and applying
\code{existsObjects\_of\_admissibleTwistDatum} with the original level-prime
theorem gives \code{ObjCP}, including its geometric object data.
This step uses the original automorphic and mass results. The preceding
integer-weight construction only uses the host and the rank bounds.

Choose the code once by \code{exists\_injective\_words}. For these same
centres and this code, \code{kit\_of\_points} makes the points into a kit,
and \code{szRow\_numAt} supplies the numerical witnesses on every point.
The relation \code{Pt.From} carries the raw host, both sets and the centres
from the specified host. Applying \code{exists\_famP\_of\_kits} to these
kits on the given \(S\) gives the claimed equality of window sets.
At each point its validity record contains the original \code{HCVal} and
the exact equality \(\eta_{\rm Sz}=2m\lceil x^{c_h}\rceil\). The integer
parameter construction and the choice of actual points can produce
different permitted thick corrections. The Fourier estimate in
Proposition~\ref{scope:reselection} is uniform over these corrections.
\end{proof}

The additional duty \code{WOK} requires a second choice of centres when
it is used. In \code{hostWW}, first form a temporary family from the
prescribed hosts and provisional \code{WinOK} centres. The original
\code{windowShareRequirementClosedForm\_proof} supplies
\code{WeightVectorExists} at every sufficiently late point. By
\code{weightWindowRequirement\_host\_of\_point}, these vectors give new
centres on the host from which that point came, zero off its top set and
satisfying both \code{WinOK} and \code{WOK}. To retain a specified host
\(H_s\) at every scale \(s\), take the scale predicate in this argument to
be existence of duty centres on \(H_s\). The temporary family has window
set \(S.\mathtt{Win}\) and supplies an index at every such scale. A second
application of the universal point construction realizes the new centres
on \(H_s\). The canonical level prime can first be installed by
\code{hostExQW\_of\_hostExW}, which preserves the raw field, the two
quadratic primes and the two finite sets.

\paragraph{Recovering the host used by the original points.}
Fix a host \(H\) used to construct the old kit. Write \(T'\) for the new
top set, \(U\subseteq H.T_2\) for the new thick set in its integer
representatives, and \(C'_2\) for the new uniform Behrend constant.
Replace just these three fields of \(H\). The range bound and absence of
three-term relations pass to \(U\), as does
\(4L_w|U|\le n_0\). The remaining changes in \code{HostGood} consist
precisely of
\[
 \tfrac12\delta_1 n\le|T'|\le2\delta_1 n,\qquad
 x^{-\varepsilon}n\le|U|,\qquad
 \frac{n_0}{L_w}\exp\!\left(-C'_2\sqrt{\log n_0\log L_w}\right)
 \le |U|
\]
and the weighted bad-orbit bound on the new top set. These are the
estimates in Proposition~\ref{scope:reselection}. The choice in
Proposition~\ref{orb:selection} makes the bad-orbit cost zero.
The canonical auxiliary level \(q_0(H)\) uses only the unchanged field
and level data. The old \(\mathrm{VVle}\) bound is therefore available
on the same field and at exactly the old rank cutoff. The coset-share
test is supplied by the choice of \(T'\).

For a signed inequality on a family assembled from the original kits,
the provenance of the point in that inequality can be retained during
assembly. The original \code{supplyChi} applied to
\code{vwin\_of\_row} supplies \(\mathrm{HostExW}(3/2,(G_\psi).\mathrm{req})\),
whose request includes both the character conditions and \(\mathrm{VVle}\).
In \code{kitS}, each scale supplies one admissible \(H\)
before the points for its translates are chosen. For every one of those
points, \code{Pt.From} identifies the ranks, the raw host, the encoded
top set and the literal integer thick set with the corresponding data
of \(H\). Its field equality and injectivity of the encoding show that
there exists a canonical-level
\code{HostGood} host with precisely these ranks, raw host and sets, and
the original \(C_2\). This property passes through the same family
assembly in the proposition. Hence a signed estimate on these kits over
complete windows supplies the old hosts required above, and all new
points and centres then follow from the original construction. The
signed estimate on every scale of those windows is still required.

For an independently specified raw field, \code{hostOfField} can also
produce a full admissible host from the page and modulus bounds,
\code{PairWin}, \code{PairGood}, \code{FieldRows} and \code{admitted}.
It chooses its own two sets. Finally an isolated \code{RealizedData}
record describes the old Fourier points, character tests and rank budget.
The standalone existence of such a record on one family and a separate
existence statement for admissible hosts on another family do not
identify their fields or windows. The conditional implication on
old kits uses their common origin at each original scale.

\begin{corollary}[One translate and one word at each scale]\label{scope:one-word}
Fix the exponent, the original cyclotomic host and full page inputs,
the uniform Fourier estimates of Proposition~\ref{scope:reselection}
and the point construction of Proposition~\ref{scope:reselected-hosts}
on one set of complete windows. At each scale retain the prescribed
thick set and choose the tops of Proposition~\ref{orb:selection}
before constructing the common points. Suppose that on these points,
at each sufficiently late scale of the windows, there exist
\(u_0\in A_0\) and one actual word \(i\) satisfying the local lattice
comparisons of Section~\ref{geom:relative-finite-index} and
\[
       (\mathcal A_{u_0}-G_{u_0})+P_{u_0i}+\ell_{u_0i}
                    \le E_x .
\]
Then both \eqref{scope:average} and \eqref{scope:missing} hold on a
family with the same windows, fields, auxiliary primes, ranks and
exponent, beyond one uniform threshold. If the right side in the test
can be replaced by \((1-\lambda)E_x\) for a fixed \(0<\lambda<1\),
the very family on which the word is tested satisfies both bounds.
\end{corollary}
\begin{proof}
Write \(B=|T|mb\), \(\delta=|T_2|/n\),
\(d_u=2D_u/(B\delta)\) and \(W=2E_x/(B\delta)\).
The word inequality and \eqref{geom:finite-cancel} with the upper
word loss give \(D_{u_0}\le E_x\). The kernel estimates in the
proof of Proposition~\ref{scope:reselection} hold on the retained
field. Corollary~\ref{orb:overlap} gives a common
\(S=S_{\psi;q,d}(x)\) and a uniform \(r_x=o(\sqrt{x})\) with
\(|d_u-S|\le r_x\) for every \(u\in A_0\).
In particular \(d_u\le d_{u_0}+2r_x\).
The original ranks give \(W\ge2\sqrt{x}\) eventually, so any
fixed margin \(\lambda W\) absorbs \(2r_x\) on this same family.

For the weak inequality, thin the thick set by the fixed factor of
Proposition~\ref{scope:reselection} and reselect its tops. The proof
of that proposition constructs the new admissible sets independently
of the old discrepancy and gives \(W'\ge2W\). Use the same input
hosts to assemble their points. With the same two primes, the new
normalized values satisfy \(d'_u\le S+r'_x\), uniformly with
\(r'_x=o(\sqrt{x})\). Consequently
\[
              d'_u\le d_{u_0}+r_x+r'_x\le W+r_x+r'_x\le W'
                        \qquad (u\in A_0)
\]
beyond one threshold. This proves the pointwise estimate, and
averaging gives its signed form on the same data. The word and the
translate may vary with \(x\). The inequality must hold after the
first top choice, at every scale of each required window.
\end{proof}
The companion \code{OneWordWindow.lean} checks the transfer to the
original subgroup and its signed mean on \code{RealizedData}, given the
single-word comparison and both Fourier bounds with their budget margin.

There is a second way to supply the fixed margin at \(u_0\).
At this translate take the word probabilities of
Corollary~\ref{rank:one-minor}. Under its per-word division, height and
variation assumptions put \(L_{u_0}=K_7+c_s I_{u_0}\), and let
\(0<\tau\le1\) be a common lower rank ratio on the proper target HN
cutoffs for these probabilities. The integral-column case of
Proposition~\ref{rank:transfer} gives the pointwise bound
\[
 D_{u_0}\le \mathcal A_{u_0}-G_{u_0}+K_0+
                 \frac{(1-\tau)K_3+L_{u_0}}{\tau}.
\]
Indeed the weighted word loss is at most \(L_{u_0}\), and the centred
maximum source slope is at most \(K_3+L_{u_0}\). On the original
sizes \(K_0\ll y^2P_0\), \(K_3+K_7\ll yP_0\). Thus at this single
translate the minor and rank budget of Corollary~\ref{rank:one-minor},
\(\mathcal A_{u_0}-G_{u_0}\le C y^2P_0\) for a uniform \(C\), and
\(c_s I_{u_0}\ll yP_0\) suffice. Indeed \(\tau\gg y^{-4+\xi}\)
and \(D_{u_0}\le C'(y^2+y^{5-\xi})P_0\) for a uniform \(C'\).
This upper bound is \(o(y^5P_0)\), while \(y^5P_0\le E_x\).
The comparison with \(S\) then transfers
the bound to the entire first family. Finding such a translate
on the actual selected host at each required scale remains open.

\subsection{Sparse scales and intervals of realized points}\label{scope:sparse-scales}
The companion \href{15-1-15-2-2026-11-22.pdf}{note on powers and
translate errors}, Sections 1--4, gives a fourth-power chord bound,
a fixed-power zero test and negative intervals for three targets.
Here the realizations on the prescribed hosts allow these analytic
estimates to be used with the actual geometric weights.

\begin{corollary}[Testing the word inequality at fourth powers]\label{scope:fourth-powers}
Retain the host and point supply on the full windows in
Corollary~\ref{scope:one-word}, the full rank page of
Corollary~\ref{pref:scale}, the uniform Behrend constant and the
discarded-twist bounds in Corollary~\ref{ar:meanhost}. In a window
\([Y,Z]\), \(Z=Y^{3/2}\), it suffices to establish the single-word
inequality of Corollary~\ref{scope:one-word} at
\[
             \{Y,Z\}\ \cup\ \{j^4:Y\le j^4\le Z,\ j\in\mathbf N\}.
\]
For arbitrarily late such windows the output family satisfies both
\eqref{scope:average} and \eqref{scope:missing} at every real scale
in those same windows, with the original exponent and coefficient.
The host realization at intermediate scales is still used.
\end{corollary}
\begin{proof}
Put \(F(x)=G_{\mathcal F_\psi}(x)\) and
\(\beta=1/2+97\varepsilon/16\). The formula in
Section~\ref{pref:section} expresses \(F(x)+x\log x\) as a linear
function plus ramps \((x/n-1)_+\) with coefficients
\(A(n)\Lambda(n)\ge0\). On each bounded interval the sum is finite.
For \(0<a<b\), \(x=(1-\theta)a+\theta b\), \(0\le\theta\le1\),
convexity of these ramps gives
\[
  F(x)\le (1-\theta)F(a)+\theta F(b)+\frac{(x-a)(b-x)}{x}.
\]
Indeed the error before bounding it is
\((1-\theta)a\log(a/x)+\theta b\log(b/x)\). Use \(\log t\le t-1\)
on both logarithms. The resulting numerator over \(x\) is the
weighted variance \((1-\theta)a^2+\theta b^2-x^2=(x-a)(b-x)\).
The companion \code{PrimeWindowInterpolation.lean} verifies this
inequality for every finite sum of nonnegative weighted ramps and
every affine term.
For adjacent test points \(a,b\), including the intervals cut by
\(Y,Z\), we have \(b-a\ll x^{3/4}\). Thus the error throughout
such an interval is at most \((b-a)^2/(4x)\ll\sqrt x\).

At a tested point \(s\), the original word bound gives a normalized
discrepancy at most \(W_s\). That discrepancy differs from the signed
four-twist sum by \(o(\sqrt{s})\). Removing the discarded twists costs
\(O_\psi(s^{1/2+\varepsilon})\). Hence
\[
    F(s)\le W_s+O_\psi(s^{1/2+\varepsilon})
         \le C\,s^\beta/\delta_s
\]
by Corollary~\ref{pref:scale}, where \(\delta_s\) is its thick density.
This step puts every test in terms of the same fixed function \(F\).
The auxiliary primes and the fields at distinct test points may vary.

At every scale \(x\) put
\(z_x=\sqrt{\log(n_x/o)\log L_{w,x}}\) in the prescribed host data.
For either adjacent test point \(s\), the Behrend bound gives
\(\delta_s^{-1}\le oL_{w,s}\exp(C_2z_s)\). The page gives
\(\log(n_x/o)=(2+8\varepsilon)\log x+O(1)\), and the width formula
in Proposition~\ref{host:behrend} gives
\(\log L_{w,x}=\log\log x+O_\varepsilon(1)\). Since
\(s=x+O(x^{3/4})\), these estimates give \(z_s/z_x\to1\) uniformly.
Thus \(\delta_s^{-1}\le\exp(Az_x)\) with a single fixed \(A>0\)
on all sufficiently late windows. The chord inequality gives
\(F(x)\le C' x^\beta\exp(Az_x)\) throughout them.

On the prescribed host at each \(x\) now use the thinning in the proof
of Corollary~\ref{scope:reselection-loss} with
\(g(x)=\exp((A+1)z_x)\). The new thick density satisfies
\(x^{-\varepsilon}\le\delta'_x\le1/(32g(x))\) eventually and its
Behrend constant is uniform. Choose the tops with the overlap of
Corollary~\ref{orb:overlap} and realize the resulting host. The
new budget is at least \(2x^\beta/\delta'_x\ge64g(x)x^\beta\).
Every new normalized discrepancy is at most
\(F(x)+O_\psi(x^{1/2+\varepsilon})+o(\sqrt{x})\), by the same
discarded-twist and overlap estimates on this host. One threshold
puts this below the new budget. This proves both inequalities on
the whole window. The number of signed word tests on it is at most
\(2+\lfloor Z^{1/4}\rfloor=O(Y^{3/8})\). Their validity on a
cofinal sequence of windows remains an arithmetic requirement.
\end{proof}

The same estimate applies to a mesh containing the two endpoints
whose adjacent points satisfy \(b-a\ll a^{(1+\beta)/2}\). Here
\(\beta<1\) since \(\varepsilon<1/52\), so \(a,b\sim x\) on such
an interval and its chord error is \(O(x^\beta)\).

If the minor and the pointwise loss bounds following
Corollary~\ref{scope:one-word} hold at the fourth powers and endpoints,
one may keep the first family with its overlap estimates at all scales.
Indeed, for \(\kappa=\min(3,\xi)>0\) these bounds give
\(D_{u_0}\le C s^{-\kappa\varepsilon}E_s\) at every tested \(s\).
The Behrend estimate gives \(\delta_s^{-1}=s^{o(1)}\) uniformly.
Thus the upper bounds for \(F(s)\) divided by \(s^\beta\) tend to
zero, since \(1/2+\varepsilon<\beta\). The fourth-power chord error
is also \(o(x^\beta)\). At every real \(x\) the same overlap and
discarded-twist bounds put all the first family's normalized values
below its own budget \(W_x\ge2x^\beta\) eventually.

\begin{proposition}[A fixed power on a cofinal set of windows]\label{scope:fixed-powers}
Use the full page, uniform host, Behrend and Fourier hypotheses of
Corollary~\ref{scope:fourth-powers}. Fix the character and the exponent,
and set \(\beta=1/2+97\varepsilon/16\). Let \(\Theta\) be the supremum
of the real parts of the nontrivial zeros in the fixed retained family.
Fix any positive integer \(r\) before the windows are selected.
Suppose the single-word test holds at every integer \(r\)-th power
in arbitrarily late complete windows, or that \eqref{scope:average}
holds at those powers on hosts satisfying its mean Fourier comparison.
Then \(\Theta\le\beta\), and both
\eqref{scope:missing} and \eqref{scope:average} hold at every
sufficiently late scale of the realized family with the overlap of
Corollary~\ref{orb:overlap}. Conversely \(\Theta\le\beta\) suffices for
these last two assertions on that family.
\end{proposition}
\begin{proof}
For a retained primitive \(\chi\) put \(b_\chi=\mathbf1_{\chi=1}\),
\(f_\chi(s)=(s-1)^{b_\chi}L(s,\chi)\), \(p=f'_\chi/f_\chi\) and
\(m_0=\operatorname{ord}_0 f_\chi\). The regularized function is entire
of order at most one, by the classical completed Hadamard product
and the reciprocal gamma factor. See
\href{https://dms.umontreal.ca/~koukoulo/documents/publications/primes.pdf}{Koukoulopoulos,
\emph{The Distribution of Prime Numbers}}, Chapters 8 and 11.
Write \(\rho\) for all its nonzero zeros, counted with multiplicity,
including the trivial ones. Subtract the logarithmic derivatives of
the genus-one product at \(1\) and \(s\). This gives
\[
 \frac{p(1)-p(s)}{s(s-1)}
    =\frac{m_0}{s^2}
       +\sum_{\rho\ne0}\frac{1}{(1-\rho)s(s-\rho)}.
\]
The left side is the Mellin transform of \(G_\chi(x)-G_\chi(1)\)
by \eqref{eq:mellin}. Since \(\sum_{\rho\ne0}|\rho|^{-3/2}<\infty\),
termwise integration on \(\Re s>1\) and uniqueness of this transform
for continuous functions give
\[
 G_\chi(x)-G_\chi(1)=m_0\log x+
             \sum_{\rho\ne0}\frac{x^\rho-1}{\rho(1-\rho)}
                    \qquad(x\ge1).
\]
The denominator also accounts for trivial zeros other than zero.
This proves \(|F(x)|\le C_\psi x^\Theta\) for \(x\ge2\), including
the endpoint exponent: the family contains the zeta function, so
\(\Theta\ge1/2\). For \(0<\Re\rho<1\) and \(|v|\le1/2\),
\[
 |(1+v)^\rho-1|\ll
   \min(1,(1+|\rho|)|v|)\le (1+|\rho|)^{1/2}|v|^{1/2}.
\]
For a negative trivial zero, the mean value theorem bounds the change
of \(t^\rho\) by \(2|\rho|(t/2)^\rho |v|\) for \(t\ge4\).
Summing the series and the change in \(m_0\log t\) yields
\[
    |F(t(1+v))-F(t)|\ll_\psi t^\Theta |v|^{1/2}.
\]

At a tested power \(s\), the single-word comparison, or the mean
Fourier formula if the mean is given, implies
\(F(s)\le C s^\beta/\delta_s\). Here \(\delta_s^{-1}=s^{o(1)}\)
uniformly. Every \(x\) in a sufficiently large complete window
is within \(O_r(x^{1-1/r})\) of an integer \(r\)-th power in that
window. At the ends choose the adjacent power towards the interior.
If \(\Theta>\beta\), choose \(0<\eta<\Theta-\beta\). The two bounds give
\[
 F(x)\le C_\eta x^{\beta+\eta}+C' x^{\Theta-1/(2r)}
\]
on the full windows. A fixed exponent strictly between the larger
of these exponents and \(\Theta\) now contradicts
Theorem~\ref{an:zero}. Thus \(\Theta\le\beta\).

On the realized hosts with the overlap estimate we have
\(d_u\le F(x)+O_\psi(x^{1/2+\varepsilon})+o(\sqrt{x})\).
The thick-set upper bound \(4 L_{w,x}|T_{2,x}|\le n_x/o\) gives
\(\delta_x^{-1}\ge4oL_{w,x}\to\infty\). Consequently
\(W_x\ge2x^\beta/\delta_x\) absorbs the fixed constant in
\(|F(x)|\le C_\psi x^\Theta\) even when \(\Theta=\beta\).
This proves the two original inequalities at the same exponent
on all sufficiently late scales where these hosts are supplied.
\end{proof}

In this last argument the mean witnesses and the output points
may belong to different families. Their Fourier comparisons both
use the fixed retained function of \(\psi\). To keep the fields and
windows of the input family, use its prescribed hosts for the
output points as in Proposition~\ref{scope:reselected-hosts}.

\begin{corollary}[Intervals of actual points for three targets]\label{scope:negative-points}
For the zeta target and for either nonprincipal primitive real
character of conductor three or four, take the common admissible
hosts supplied by the original construction with the page and Fourier
bounds above. Select and realize the first family with the overlap
of Corollary~\ref{orb:overlap}. For every fixed \(0<\nu<1\), every
sufficiently late complete window of this family contains an interval
\[
       [t-t^{1-\nu},\ t+t^{1-\nu}]
\]
on which both \eqref{scope:average} and \eqref{scope:missing} hold for
its actual points. This conclusion requires no signed-word estimate.
\end{corollary}
\begin{proof}
Each of these retained families has a nontrivial zero and has none on
the real interval \((0,1)\). For zeta use the alternating series,
and for the other two characters pair each positive term with the
next negative one in its Dirichlet series. Their retained families
contain zeta. See also the
\href{https://dlmf.nist.gov/25.10}{distribution of zeta zeros}.
Fix \(\max(0,\Theta-\nu/2)<\theta<\Theta\). Every sufficiently late
power window of fixed relative logarithmic width contains some \(t\)
with \(F(t)<-t^\theta\). Otherwise the proof of Theorem~\ref{an:zero}
applied to \(-F\) on arbitrarily late windows gives a contradiction:
only the nonreal case of Lemma~\ref{an:isolate} can occur, and its
lower bound uses the absolute value of the nonzero residue.
Take \(Y^{7/6}\le t\le Y^{4/3}\) within the prescribed window.
The local estimate for \(F\) in the preceding proof bounds its
change on the stated interval by \(O_\psi(t^{\Theta-\nu/2})\).
Hence \(F(x)<0\) throughout it, and both endpoints lie in
\([Y,Y^{3/2}]\) for large \(Y\).

At every such \(x\) the overlap estimate for the actual weights and
the discarded-twist bound give
\[
 \max\!\left(0,\max_{u\in A_0}\frac{D_u}{E_x}\right)
       \ll_{\psi,\varepsilon} x^{-81\varepsilon/16},
\]
using \(W_x\ge 2x^{1/2+97\varepsilon/16}\). A single threshold suffices
for all translates. The original host and point supply is used on its
own complete windows. The interval occupies a vanishing fraction of
the scale \(t\) for fixed \(\nu>0\).
\end{proof}
